\chapter{Conclusão}
\label{ch:conclusao}

\pendente{retomar o problema de pesquisa e responder em um parágrafo ao objetivo geral (Seção~\ref{sec:obj-geral})}

\section{Atendimento aos objetivos específicos}
\label{sec:conc-objetivos}

\pendente{para cada objetivo específico da Seção~\ref{sec:obj-especificos}, indicar em que seção ele foi atendido e qual foi o principal achado}

\section{Contribuições}
\label{sec:conc-contribuicoes}

\pendente{resumir os três artefatos (Tabelas~\ref{tab:requisitos}, \ref{tab:problemas-ux} e \ref{tab:taxonomia-arquitetura}) e para quem cada um é útil: gestores, fornecedores de soluções e pesquisadores}

\section{Limitações}
\label{sec:conc-limitacoes}

Os resultados devem ser lidos considerando as limitações a seguir.

\begin{itemize}
    \item \textbf{Amostra por conveniência.} O questionário foi divulgado na rede pessoal e acadêmica do autor, com repasse entre contatos. A amostra concentra respondentes jovens adultos e com ensino superior, e tem poucos respondentes com deficiência, o que impede generalizar os resultados para todos os frequentadores de Porto Alegre e limita a discussão sobre acessibilidade.
    \item \textbf{Janela de coleta curta e sem piloto.} O questionário ficou aberto por poucos dias e não passou por aplicação piloto com respondentes externos.
    \item \textbf{Construtos com dois itens.} Cada construto da UTAUT2 foi medido por duas afirmações. Em alguns construtos, os dois itens se mostraram pouco correlacionados, o que indica que medem aspectos diferentes; nesses casos, a análise foi feita item a item.
    \item \textbf{Número de casos.} O estudo de caso múltiplo envolveu quatro estabelecimentos, um por perfil, e um único fornecedor, o que permite contraste analítico, mas não comparação entre casos do mesmo perfil nem entre fornecedores. O perfil D foi incorporado durante a coleta, e não no desenho original.
    \item \textbf{Proximidade entre pesquisador e entrevistados.} Parte dos entrevistados foi acessada pela rede pessoal do autor, e em alguns momentos o entrevistador expressou concordância com o que ouvia. Para reduzir esse efeito, a análise privilegia as falas espontâneas, anteriores a qualquer comentário do entrevistador.
    \item \textbf{Gravações incompletas.} Em ao menos uma entrevista, o início da conversa não foi transcrito e a gravação disponível ficou incompleta; o contexto ausente foi registrado pelo pesquisador logo após a entrevista.
    \item \textbf{Dados declarados.} Tanto o questionário quanto as entrevistas registram percepções e relatos, sem observação direta do uso nem medição do desempenho das soluções.
\end{itemize}

\pendente{revisar as limitações com a base final e acrescentar as que surgirem nas entrevistas}

\section{Trabalhos futuros}
\label{sec:conc-futuros}

A partir dos resultados e das limitações, sugerem-se as seguintes continuações:

\begin{itemize}
    \item avaliar interfaces de cardápio digital por meio de testes de usabilidade com usuários, incluindo pessoas idosas e pessoas com deficiência visual ou de leitura;
    \item validar uma escala com mais itens por construto, em amostra maior e mais diversa, permitindo análise fatorial e testes de hipótese;
    \item comparar cidades com e sem exigência legal de cardápio físico, para observar o efeito da regulação sobre a experiência dos clientes;
    \item prototipar uma solução de cardápio digital orientada pelo catálogo de requisitos deste trabalho e avaliá-la com gestores e clientes.
\end{itemize}
